\documentclass[../main.tex]{subfiles}
\begin{document}


\section*{第一部分提要}\label{section:part1-outline}

\cref{chapter:ch01}给出抽象测度论的一般结果，内容包括外测度与抽象容量、Carathéodory 构造、Hausdorff 测度与容量，以及分解与向量测度的一般结论。

在\cref{chapter:ch02}中，我们先回顾标量函数可测性与积分的经典结果，然后引入两个对本书展开至关重要的主题：Young 测度与 Riesz 表示定理（包括正泛函与有界泛函两种情形）。随后讨论 Lebesgue 空间中各种收敛概念，并指出一致可积性的重要作用（Vitali 定理、Dunford--Pettis 定理）。最后一节处理微分理论的一个主要结果，即 Radon--Nikodým 定理，以及 Euclid 空间 $\mathbb{R}^N$ 上 Radon 测度的微分。

\cref{chapter:ch03}首先介绍 Sobolev 与 Bessel 位势空间的若干材料，然后回顾有界变差函数与 Sobolev 函数的结果。在此基础上，我们详细讨论 $\mathbb{R}^N$ 子集上的各种容量概念（Bessel、Riesz、Sobolev 与 Hausdorff 容量）及其相互关系。Bessel 容量与 Riesz 容量的结果由与核相伴的 Meyer 容量理论导出。关于 Sobolev 容量集中或弥散的测度也将予以讨论并刻画。

\cref{chapter:ch04}处理向量积分。在引入向量函数可测性的若干概念并讨论其相互关系（特别是 Pettis 定理）之后，我们给出 Bochner 积分以及较弱积分概念的相关结果。随后引入向量 Lebesgue 空间 $L^p(X;Y)$、$L^p_w(X;Y)$ 与 $L^p_{w^*}(X;Y^*)$（其中 $p\in[1,\infty]$，$X$、$Y$ 为 Banach 空间），并证明完备性与可分性结果。然后讨论这些空间的对偶理论（特别是当 $Y^*$ 或 $Y$ 可分时 $L^p_{w^*}(X;Y^*)$ 的对偶），指出所谓 Radon--Nikodým 性质的核心地位，等价地说，即从 $L^1(X)$ 到 $Y$ 的线性连续算子的 Riesz 可表示性。我们所着眼的应用是 $Y=C_0(Z)$、$Y^*=\mathscr{R}_f(Z)$ 的情形，其中 $\mathscr{R}_f(Z)$ 为某度量空间 $Z$ 上有限 Radon 测度的 Banach 空间；因为空间 $L^\infty_{w^*}(0,T;\mathscr{R}_f(\Omega))$（$T>0$，$\Omega\subseteq\mathbb{R}^N$）正是本书第二部分所处理发展问题的自然框架。该章最后两节分别处理 $Y$ 为实函数空间时的向量 Lebesgue 空间，以及向量 Sobolev 空间。

\cref{chapter:ch05}致力于有限 Radon 测度序列的收敛性。Banach 空间 $\mathscr{R}_f(X)$ 中的有界性保证弱$^*$拓扑下的紧性（由 Riesz 定理与 Banach--Alaoglu 定理），但在更强的窄拓扑下的紧性则同时要求有界性与胎紧性，这一概念在此处起着核心作用。我们给出窄收敛的判别准则（由 Portmanteau 定理给出，或利用 Prokhorov 距离）以及胎紧性的判别准则。

随后证明 Hausdorff 空间乘积上有限 Radon 测度的离散分解定理。将离散分解定理应用于与可测函数 $u_j:X\to\mathbb{R}$（$j\in\mathbb{N}$）相伴的 Young 测度序列 $\{\nu_{u_j}\}$，可以建立序列 $\{u_j\}$ 在 Lebesgue 空间中的有界性与 $\{\nu_{u_j}\}$ 关于窄拓扑的（相对）紧性之间的关系。特别地，若 $\{u_j\}$ 在 $L^1(X)$ 中有界，则相伴 Young 测度序列 $\{\nu_{u_j}\}$ 在窄意义下（可能取子列）收敛于某个 Young 测度 $\nu$。另一方面，由 Banach--Alaoglu 定理，存在 $\{u_{j_k}\}\subseteq\{u_j\}$ 与 $\sigma\in\mathscr{R}_f(X)$，使得 $u_{j_k}\stackrel{*}{\rightharpoonup}\sigma$ 在 $\mathscr{R}_f(X)$ 中。Biting 引理（亦将证明）给出与这种收敛相关的集中现象的精确描述，并给出极限测度 $\sigma$ 与上述 Young 测度 $\nu$ 之间的关系。

\end{document}
